\subsection{学习表示与利用表示结构}
OPQ学习旋转和子空间码本以降低失真\cite{opq}；JHQ的随机旋转则支持解析primary模型。
通用学习码本未必保留$16+16$精确分解所需的Cartesian结构。
ScaNN按最大内积排序重要性加权量化误差\cite{scann}，说明MSE不是唯一目标；
本文评估平方L2下近邻保留与GPU执行成本。

\subsection{小查找表已有充分先行工作}
PQ Fast Scan将表重组为适于CPU SIMD的寄存器查找\cite{fastscan}，
Quicker ADC扩展到不规则量化器和split tables\cite{quickeradc}。
JHQ利用内部Cartesian结构实现GPU因子化；FP32结合顺序及截断并列仍需数值检验。

RaBitQ用随机化与紧凑符号信息建立特定距离估计器和误差界\cite{rabitq}；
多bit扩展以存储换精度\cite{rabitq_multi}。GPU IVF-RaBitQ进一步融合GPU建库、
评分与搜索\cite{ivfrabitq}，是直接竞争基线。其理论保证不能因表示相似而转移给JHQ；
JHQ以主符号筛选候选，再显式读取残差层。

\subsection{精排工作量与自适应停止}
Faiss的\texttt{k\_factor}与cuVS refinement已提供候选扩张控制
\cite{faiss_refine,cuvs_refinement}。Faiss autotuning用GT探索recall--时间取舍
\cite{faiss_autotune}；学习式终止预测搜索充分性\cite{li_termination}，
DARTH以声明式recall目标组织早停\cite{darth}。
第6章以无标签参考一致性为整个工作负载选择一个倍率，无需预测器训练，
但继承参考自身的误差。
